% \section{Mapping of heart beat to music}
% den abschnitt nocheinmal prüfen 

% RPPG Techniques provide the functionality to measure the heart beat remotely. These measurements then need to be converted to music.

% There has been some research on passive music listening without any biofeedback mechanism. According to a meta-analysis by  Pelletier(2004), music consistently assists in lowering stress levels. The magnitude of this response is affected by several factors. Patients with prior musical experience, such as playing an instrument, tend to lower their stress levels a lot more.
% Furthermore, Pelletier notes that indivudal therapy generally results in better relaxation than group session. Although, group sessions are used much more often because they are less expensive. A huge obstacle is that the devices that are used for data collection often cause discomfort. This can lead to false data regarding the relaxation process. Finally, the study suggests that music based on scientific research (e.g., specific rhythmic or frequency patterns) is more effective for stress reduction that music chosen soley based on personal preference \cite{pelletier2004effect}.
% In another study, Linnemann et al. examined how music affects stress in daily life rather than in a laboratory. They came to the conclusion that listening to music can effectively reduce stress levels. Nevertheless, the largest effect is reached when subjects listen to music with then intention of decreasing stress levels. This suggests that the reason why someons listens to music is very important for the actual stress-reducing effect \cite{linnemann2015music}.

% These results are also supported by another study \cite{thoma2011listening}. They confirm that listening to music alone does not automatically decrease stress levels. Instead,  the reason of listening to the music is the most important factor. For example, if someone listens to music to cope with loneliness, it might effectively decrease their feeling of isolation. Nevertheless, if the goal is relaxation, the music only works well if the listener has the specific intention to relax \cite{thoma2011listening}.

% The study of Bernardi et al. (2006) confirms that the results do not depend on the music preferation of the user rather on the BPM and rhythm structure of the music. Additionally, the study found no evidence of habituation, meaning the users do not simply 'get used' to the music over time. In Contrast to Linnemann et al. (2015), who emphasized the importance of the listener's intention, Bernadi concludes that music can be perceived as a passive form of meditaiton. consequently, this means that the user does not need an intend to reduce stress for measurable results \cite{bernardi2006cardiovascular}.

% Passive music listening without any biofeedback system has mixed results. To convert the heartbeat to music, a process called mapping is necessary, that takes the input data and transforms it into music. There are two methods Continous mapping and diskrete mapping. Continous mapping describes the process of constantly changing the music according to the data. During the discrete mapping process certain aspects of the music are changed when a certain threshold is reached \cite{bergstrom2014using}.

\section{Biofeedback-Driven Regeneration}

As explained previously the heart rate is converted to a piece of music. It is now crucial to understand how the music can help the human body to regenerate. According to Lehrer and Gevirtz (2014), the efficiency of biofeedback relies on the synchronization of the cardiovascular and respiration systems \cite{lehrer2014heart}.

The human cardiovascular system has a certain resonance frequency at which point it exhibits the best effects on the human body. If this state is reached, has maximum internal order and physiological efficiency. This frequency is reached at approximately 0.1 Hz. This is equivalent to roughly six breaths per minute. By breathing at the resonance frequency, heart rate and blood pressure enter a harmonic state. Furthermore, gas exchange in the lungs are also optimised at this state \cite{lehrer2014heart}.

In addition, the activity of the vagus nerve is stimulated. The vagus nerve is responsible for the communication between the heart and the central nervous system. This makes it possible for the body to lower the heart rate more effectively \cite{lehrer2014heart}.

Moreover, biofeedback creates a meditative effect that changes the user's state of consciousness. The user has to focus his thoughts and his energy on the biofeedback signal. Thus, it the brain cannot focus on anxious thoughts. This shift in mental focus leads to a state of heightened awareness, which is explored in the following section on consciousness and mindfulness \cite{lehrer2014heart}.